\chapter*{Tóm tắt}
\addcontentsline{toc}{chapter}{Tóm tắt}
\markright{Tóm tắt}

Tinh chỉnh (fine-tuning) một mô hình ngôn ngữ huấn luyện trước (pre-trained) cho từng tác vụ mang lại hiệu năng tốt trên tác vụ đó, nhưng mỗi tác vụ mới lại thêm một mô hình hoàn chỉnh cần huấn luyện, lưu trữ và vận hành, đồng thời những gì học được cho một tác vụ không đến được các tác vụ khác. Nghiên cứu này xem xét liệu nhiều mô hình đã tinh chỉnh có thể trở thành một mô hình đa nhiệm duy nhất với trọng số thích ứng theo đầu vào hay không. Chúng tôi áp dụng Weight-Ensembling Mixture of Experts (WEMoE), vốn được đánh giá trên các vision Transformer, cho BERT và năm tác vụ nhị phân của bộ đánh giá GLUE: MRPC, RTE, CoLA, SST-2 và QNLI. Các trọng số ngoài khối MLP được hợp nhất tĩnh bằng Task Arithmetic, còn trọng số MLP của mỗi lớp encoder được dựng lại cho từng đầu vào bằng trọng số huấn luyện trước cộng với tổng các véc-tơ tác vụ đóng băng, theo hệ số do bộ định tuyến (router) quyết định. Chỉ các bộ định tuyến được huấn luyện, bằng Few-Shot Fine-Tuning (FSFT) trên các mẫu có nhãn mà các chuyên gia (expert) chưa từng thấy. Mọi lựa chọn mô hình đều dùng tập phát triển (dev), và điểm được báo cáo trên các tập kiểm định chính thức. Với 256 mẫu mỗi tác vụ, WEMoE đạt điểm trung bình $69.9\pm0.1$, so với 57.4 của Task Arithmetic, 78.2 của các chuyên gia và 45.8 của bộ dự đoán lớp đa số, tương ứng 59.1\% mức tăng giữa lớp đa số và các chuyên gia. Các hệ số không phụ thuộc đầu vào chỉ đạt $58.4\pm0.0$. Chỉ nâng cấp lớp đầu tiên của MLP cho điểm trung bình tương tự ($69.1\pm0.1$) với 47.2\% dung lượng lưu trữ của năm chuyên gia thay vì 73.1\%, nhưng hai biến thể yếu ở các tác vụ khác nhau. WEMoE không bằng chuyên gia ở tác vụ nào và chậm hơn nhiều lần khi chạy theo batch. Các hạn chế gồm một bộ chuyên gia duy nhất, tập kiểm định nhỏ ở hai tác vụ, và dữ liệu dev vượt quá ngân sách few-shot.

\vspace{3mm}
\noindent\textbf{Từ khóa:} hợp nhất mô hình; weight ensembling; Mixture of Experts; BERT; học đa nhiệm; few-shot fine-tuning.
